\section{两类错误与功效共同定义检验}
\label{sec:11-Mathematical-Statistics/05-Hypothesis-Testing/01}

实验平台上挂着一条结论：新排序模型对照旧模型，转化率 \(5.02\%\) 对 \(4.97\%\)，\(p=0.31\)，判定``无显著差异''，需求关闭。三个月后另一个团队把同一份改动重新上线，这次每组跑了六倍的量，测出 \(0.24\) 个百分点的提升，\(p=0.004\)。

第一次实验没有测出效应，它同样没有测出``没有效应''。按当时每组两万人的样本量算，若真实提升就是 \(0.24\) 个百分点，那次实验交出``不显著''的概率约为 \(0.79\)。它从设计的那一刻起，就有近八成的概率得到后来被推翻的那个结论。做与不做，差别不大。

要看清这件事，得先把检验当成它实际上的那个东西：一条把数据映射到二元行动的规则。规则会出错，而且有两种错法，评价一条规则必须把两种错法一起看。

\subsection{一条规则的两种错法}

假设检验先把参数空间切成两块：

\[
H_0:\theta\in\Theta_0,
\qquad
H_1:\theta\in\Theta_1 .
\]

再指定一个拒绝域 \(R\)，数据落进去就宣布``拒绝 \(H_0\)''，否则``不拒绝 \(H_0\)''。注意后一句的措辞：证据不足以拒绝，与证明原假设成立，是两回事。这条措辞上的谨慎在本单元里会反复用到。

真实参数落在 \(\Theta_0\) 里而规则拒绝了它，是第一类错误，俗称弃真。检验的水平取所有原假设参数上的最坏情况：

\[
\sup_{\theta\in\Theta_0} P_\theta(X\in R)\le\alpha .
\]

真实参数落在 \(\Theta_1\) 里而规则没有拒绝，是第二类错误，俗称取伪，概率记作

\[
\beta(\theta)=P_\theta(X\notin R),\qquad \theta\in\Theta_1 ,
\]

功效则是它的补：\(\pi(\theta)=1-\beta(\theta)=P_\theta(X\in R)\)。

两个量的写法不对称，这处不对称藏着实践中的大半麻烦。\(\alpha\) 是一个数，用上确界一次性锁死；\(\beta\) 是一个函数，效应越小它越大，趋近于 \(1-\alpha\)。所以``这个检验的功效是 \(80\%\)''单独说出来没有内容，必须补齐针对多大的效应、多大的噪声、多少样本。汇报里凡是出现不带效应量的功效数字，基本可以断定没人算过。

一个例子把两侧接上。设 \(X_i\sim N(\mu,\sigma^2)\)，\(\sigma\) 已知，检验 \(H_0:\mu=\mu_0\) 对 \(H_1:\mu>\mu_0\)，统计量取 \(Z=(\bar X-\mu_0)/(\sigma/\sqrt n)\)，拒绝域取 \(Z>z_{1-\alpha}\)。在 \(H_0\) 下 \(Z\sim N(0,1)\)，所以第一类错误恰为 \(\alpha\)。若真实均值是 \(\mu_1=\mu_0+\Delta\)，功效为

\[
\begin{aligned}
\pi(\mu_1)
&=P_{\mu_1}\!\left(\bar X>\mu_0+z_{1-\alpha}\frac{\sigma}{\sqrt n}\right)\\
&=P_{\mu_1}\!\left(\frac{\bar X-\mu_1}{\sigma/\sqrt n}>z_{1-\alpha}-\frac{\sqrt n\,\Delta}{\sigma}\right)\\
&=\Phi\!\left(\frac{\sqrt n\,\Delta}{\sigma}-z_{1-\alpha}\right).
\end{aligned}
\]

骨架三步：把拒绝域从 \(Z\) 的语言翻回 \(\bar X\) 的语言；在真实分布下重新标准化，中心从 \(\mu_0\) 挪到 \(\mu_1\)，多出的位移就是 \(\sqrt n\,\Delta/\sigma\)；用对称性把上尾概率写成 \(\Phi\)。条件对账：第一步只用了拒绝域的单调形式；第二步用到 \(\sigma\) 已知且不随均值变化，方差未知时这里换成 \(t\) 分布且带非中心参数，量级结论不变；第三步用到正态的对称性。整个式子只依赖一个组合量 \(\sqrt n\,\Delta/\sigma\)，样本量、效应、噪声三者从此绑成一根绳。

\subsection{临界值只是在两类错误之间搬运}

拒绝域的位置一移，两个概率同时改变，而且方向相反。把阈值 \(c\) 往右推，\(H_0\) 的右尾变小，第一类错误下降；同一条线也把 \(H_1\) 分布的更多质量留在了不拒绝的一侧，第二类错误上升。

\begin{center}
\begin{tikzpicture}[x=0.86cm,y=1.5cm,font=\small,>=stealth]
  \fill[black!18]
    plot[domain=-0.4:1.75,samples=50] (\x,{1.55*exp(-0.5*(\x-3)*(\x-3))})
    -- (1.75,0) -- (-0.4,0) -- cycle;
  \fill[black!55]
    plot[domain=1.75:3.6,samples=50] (\x,{1.55*exp(-0.5*\x*\x)})
    -- (3.6,0) -- (1.75,0) -- cycle;
  \draw[black!75,thick] plot[domain=-3.4:3.6,samples=90] (\x,{1.55*exp(-0.5*\x*\x)});
  \draw[black!75,thick,densely dashed] plot[domain=-0.4:6.4,samples=90] (\x,{1.55*exp(-0.5*(\x-3)*(\x-3))});
  \draw[black!60] (-3.6,0) -- (6.6,0);
  \draw[black!85] (1.75,-0.06) -- (1.75,1.72);
  \node[font=\footnotesize,black!80] at (0,1.78) {\(H_0\) 下的抽样分布};
  \node[font=\footnotesize,black!80] at (3.9,1.78) {\(H_1\) 下的抽样分布};
  \node[font=\footnotesize,black!85] at (1.75,-0.30) {临界值 \(c\)};
  \node[font=\footnotesize,black!85] at (2.55,0.22) {\(\alpha\)};
  \node[font=\footnotesize,black!85] at (0.95,0.20) {\(\beta\)};
  \draw[->,black!70] (1.9,-0.62) -- (3.4,-0.62);
  \node[font=\footnotesize,black!70,anchor=west] at (3.55,-0.62) {\(c\) 右移：\(\alpha\) 减小，\(\beta\) 增大};
  \draw[->,black!70] (1.6,-0.62) -- (0.1,-0.62);
  \node[font=\footnotesize,black!70,anchor=east] at (-0.05,-0.62) {\(c\) 左移：\(\alpha\) 增大，\(\beta\) 减小};
\end{tikzpicture}

\emph{两条抽样分布重叠，临界值把它们切开：右边深色是 \(H_0\) 被误拒的概率，左边浅色是 \(H_1\) 被漏掉的概率。}
\end{center}

图上唯一能自由选择的是那条竖线的位置，而两块阴影的和不会因为挪它而变小多少——挪走的部分从一块转移到另一块。要让两块同时变小，只能改动图本身：把两条曲线推得更开（增大效应，通常做不到），把它们压得更窄（减小噪声，或者增加样本量），或者换一个统计量让同样的数据分得更开（本单元后面讲似然比的那一节回答这个问题）。样本量 \(n\) 进入的正是曲线宽度：抽样分布的标准差是 \(\sigma/\sqrt n\)，\(n\) 变四倍，宽度减半，重叠部分显著缩小。

顺带说清另一个常被误用的选择。若只有``提高''才触发行动，可以预先取单侧备择，同样的 \(\alpha\) 全部押在一侧，针对该方向的功效更高。前提是``预先''。先看估计的方向再决定用哪一侧，等于在两个检验里挑对自己有利的那个，实际的第一类错误接近 \(2\alpha\) 而非宣称的 \(\alpha\)。

\subsection{四个量里知道三个就能算第四个}

把上面的功效公式反解，实验设计的全部算术就出来了。要求 \(\pi(\mu_1)\ge 1-\beta\)，即 \(\sqrt n\,\Delta/\sigma-z_{1-\alpha}\ge z_{1-\beta}\)，于是

\[
n\ \ge\ \frac{(z_{1-\alpha}+z_{1-\beta})^2\,\sigma^2}{\Delta^2}.
\]

样本量、效应量、显著性水平、功效，四个量被一个等式拴在一起，给定其中三个，第四个就定了。实验设计能做的事情，到此为止。

代进一次真实的 A/B 实验。基线转化率 \(p=5\%\)，两组等分，希望能发现 \(0.25\) 个百分点的绝对提升，双侧 \(\alpha=0.05\)（\(z=1.96\)），目标功效 \(0.8\)（\(z=0.84\)）。两组比例之差的方差是 \(2p(1-p)/n\)，代入得每组约需 \(1.2\times 10^{5}\) 人。若手上每组只有两万人，把式子反过来解 \(\Delta\)，这份样本能可靠发现的最小效应是 \(0.61\) 个百分点，相对提升 \(12\%\)。而在真实效应为 \(0.25\) 个百分点时，功效只有 \(0.21\)——开头那次被关掉的实验，就站在这个数上。

于是有三种把四元关系用起来的方式，对应三个不同的时刻。开工前解 \(n\)，这是样本量估算；样本量被排期或流量卡死时解 \(\Delta\)，得到``这个实验能看见多大的效应''，若它比业务上任何值得上线的改动都大，那么这个实验不必做；实验做完得到不显著的结果时，回头算此设计对最小有意义效应的功效，用来判断这条``不显著''值多少钱。三件事用的是同一个等式。

方差 \(\sigma^2\) 一般得靠历史数据估计，估偏了样本量就偏。指标越是长尾（人均时长、人均消费、GMV），方差越大，同样的相对提升需要的样本越多，这也是为什么转化率类指标常能在一周内跑出结论，而收入类指标往往需要几周。聚类抽样（按用户而非按会话随机化）还会进一步抬高有效方差，若忽略它，算出来的样本量会系统性偏小。

\subsection{功效不足的实验做了等于没做，甚至更糟}

低功效的第一重后果显而易见：大部分真实的改进会被判成不显著。第二重后果没那么显眼，却更伤人。

在低功效的设计里，能越过显著线的估计值必须足够大。开头那个例子里，标准误约 \(0.22\) 个百分点，要显著就得测出至少 \(1.96\times 0.22\approx 0.43\) 个百分点，而真实效应只有 \(0.25\)。这意味着：凡是这次实验报出显著的场合，报出的效应至少是真值的 \(1.7\) 倍。低功效不只让人错过真效应，它还保证了侥幸捞到的那些``成功''被系统性地夸大。后续按这个数字排期、算 ROI、定下一季度目标，偏差会一路传下去。这个现象在文献里叫夸大比（type M error），它是低功效的必然产物，与数据质量无关。

一份可以照做的清单。开工前把最小有意义效应写进文档，它是业务判断而不是统计判断——低于这个提升就不值得上线，那么也不值得为它设计实验。据此算样本量与运行天数，写进预注册，并且约定中途不看结论（为什么中途看会破坏 \(\alpha\)，下一节会算）。报告时给出效应估计与置信区间，让读者看见不确定性的宽度，这正是\hyperref[sec:11-Mathematical-Statistics/04-Interval-Estimation/01]{``置信度属于程序，而非这次参数''}里区间的用法。得到不显著结论时，同时报出该设计能检测的最小效应，这一句话可以把``没有效应''和``看不见效应''分开。

最后一条边界。若目标本来就是证明两者差别小到可忽略——新架构上线不掉指标、模型压缩后精度不降——普通检验的``未拒绝''给不了这个结论，因为它对功效不设任何要求。此时应改用等价性检验：先定一个实际容忍界 \(\delta\)，再去拒绝``差异大于 \(\delta\)''。原假设与备择假设互换位置，举证责任才落到正确的一边。

两类错误的框架到这里是完整的：一条规则，两个概率，四个量的联动。剩下的问题是，实践里几乎没有人直接报告拒绝域和 \(\alpha\)，报告里只有一个 \(p\) 值。它和这套框架是什么关系，又为什么被误读得如此普遍，是下一节的事。
